\begin{definition}[The local dimension on 12:120]
\label{c:local-dimension}
For a finite-type morphism $f:X\to Y$ of universally catenary noetherian schemes we use
\[
 d_y(f)=\sup_{\substack{x\in X\\ f(x)\text{ generalizes }y}}
 \left\{\operatorname{trdeg}_{k(f(x))}k(x)
       +\dim\overline{\{f(x)\}}^{\,\Spec\mathcal O_{Y,y}}\right\}.
\]
The closure in the second summand is taken in the local base. This is the dimension function, normalized to be zero at $y$, pulled back to $X$. It is the formula visible in the supplied transcription, whose notation for the local spectrum is uncertain. It measures local source dimension, not maximum fibre dimension: for $\Spec K\hookrightarrow\Spec R$ with $R$ a discrete valuation ring and $y$ its closed point it equals $1$, although all nonempty fibres have dimension zero. After an equidimensional localization of the source and a compactification, its maximum is the local dimension of that compactification. Regular alterations preserve this maximum. The dimension-function convention and the corresponding local cohomological interpretation are made precise in \citet[\S1.1.1 and Corollary 1.2.2]{c_pillonistroh2014}.
\end{definition}

\begin{lemma}[Proper cohomology with simultaneous geometric bounds]
\label{c:proper-level}
\status{Proved}
Let $Z$ be proper of dimension at most $d$ over a field $k$, and let $a\geq0$. After the finite field extensions allowed in the admissibility definition, $H^a(Z_{\bar k},\Q_\ell)$ has a finite filtration by subquotients of special objects $V_r(-t)$ satisfying
\[
 t\geq\max(0,a-d),\qquad r\geq0,\qquad r+2t\leq a.
\]
In particular its weights are at most $a$ and its geometric levels are at most $\min(a,2d-a)$. A negative upper level bound means the object is zero. These are geometric generation assertions, not deductions from a Hodge polygon.
\end{lemma}

\begin{proof}
For a smooth projective variety of dimension $d$, hard Lefschetz identifies $H^a$ with $H^{2d-a}(-(a-d))$ when $a>d$; for $a\leq d$ use $H^a$ itself. The input is \citet[Theorem 4.1.1]{c_deligneweil1980}. If a special object of degree $r$ is required to be presented by a variety of dimension exactly $r$, take a smooth complete intersection of dimension $r$ in the smooth projective variety. Weak Lefschetz injects $H^r$ into its middle cohomology. The restriction pairing is nondegenerate by hard Lefschetz on the original variety, so this injection even has a Galois-equivariant left inverse. A finite extension of $k$ suffices to choose the hyperplanes. Thus both conventions of specialness give the same objects over a field.

For an integral proper $Z$, \citet[Theorem 7.3 and Corollary 7.4]{c_dejong1996}, applied over an algebraic closure and then descended to a finite extension, give a smooth projective $W$, a finite group $G$ acting on it, and a proper map $V=W/G\to Z$ which is a universal homeomorphism over a dense open $U\subset Z$. The construction and the finite group action descend because they involve finitely many algebraic coefficients. Passing through finite purely inseparable extensions does not change the \(\ell\)-adic topos. Let $S=Z\setminus U$ and $E=V\times_ZS$. Both have dimension at most $d-1$. Proper base change identifies
\[
 H^a(V,\Q_\ell)=H^a(W,\Q_\ell)^G:
\]
the finite quotient has no higher direct images of the constant sheaf, and the invariant functions on every geometric orbit are the constants. Averaging is exact over $\Q_\ell$, even if $\ell$ divides $|G|$.

The localization triangles for $U\subset Z$ and its inverse image, together with proper base change, give the abstract blowup sequence
\[
 H^{a-1}(E)\longrightarrow H^a(Z)
 \longrightarrow H^a(V)\oplus H^a(S).
\]
The induction hypothesis for $E$ gives $t\geq\max(0,(a-1)-(d-1))$ and weight at most $a-1$; for $S$ it gives an even stronger lower bound on $t$. The smooth-projective assertion for $W$ handles $V$. Subquotients and extensions preserve all three inequalities. Reducible proper schemes follow by the same localization argument for their components; nilpotents do not change the \(\ell\)-adic topos. This completes the dimension induction.

Finally finite-extension descent does not enlarge the bounds. For a finite separable extension $k'/k$, a representation is a direct summand of the induction of its restriction, by the trace divided by $[k':k]$. A smooth projective $k'$-scheme viewed as a $k$-scheme is still smooth projective of the same dimension, and realizes that induced cohomology. Purely inseparable extensions have already been handled. This justifies the descent used above.
\end{proof}

\begin{theorem}[Corrected bounds and the margin question]
\label{c:general-level-bound}
\status{Proved in the stated arithmetic finite-type setting}
Let $f:X\to Y$ be separated of finite type, with $X$ regular and $X,Y$ of finite type over $\Z[1/\ell]$ or over a field of characteristic different from $\ell$. Fix $y\in Y$ and put $n=d_y(f)$. Let $F$ be strictly special of weight $j$, in the dimension-normalized sense of the insertion on 12:111: it is a direct summand of $R^jg_*\Q_\ell$ for a smooth projective $g$ of relative dimension $j$. Set $\epsilon=1$ if no component of local dimension $n$ maps to $y$, and set $\epsilon=0$ otherwise. Put
\[
 L=\max(0,i-n+\epsilon),\qquad U=\min(i,n)+j.
\]
The stalk $(R^if_*F)_{\bar y}$ is effectively admissible, with a filtration into pieces of weights $w$ and levels bounded simultaneously by
\begin{equation}
\label{c:weight-level-polygon}
 2L\leq w\leq2U,
 \qquad \operatorname{level}\leq\min\{w-2L,\,2U-w\}.
\end{equation}
Consequently
\[
 \operatorname{weight}\leq2\min(i,n)+2j,
 \qquad
 \operatorname{level}\leq
 \begin{cases}
 \min(i,2n-i)+j,&\epsilon=0,\\
 \min(i,2n-1-i)+j,&\epsilon=1.
 \end{cases}
\]
For an open immersion at a boundary point, $\epsilon=1$. The condition that $f$ have relative dimension at most $n-1$ near $y$ also implies $\epsilon=1$. Hence the margin's replacement by $n-1$ is valid with its stated extra condition, since $\min(i,2n-1-i)\leq\min(i,n-1)$. Without that condition the corrected general bound is at most $\min(i,n)+j$, and Proposition~\ref{c:level-counterexample} prevents replacing $n$ by $n-1$ universally.
\end{theorem}

\begin{proof}
We first take $F=\Q_\ell$ and $j=0$. Work locally on $Y$ around $y$, compactify $X$ over $Y$, and retain only the closure of $X$. Applied to that compactification and its boundary, the alteration theorems of \citet[Theorem 4.1 for varieties over fields; Theorem 8.2 for arithmetic schemes]{c_dejong1996} give a proper regular alteration $h:\widetilde X\to Y$ with a strict normal-crossing boundary $D$, and an open $v:X'\hookrightarrow\widetilde X$ above $X$. The induced $\pi:X'\to X$ is proper and generically finite. The construction can be made component by component. Localization to $\Spec\mathcal O_{Y,y}$ preserves regularity, boundary codimensions, and the local dimension function.

There are maps
\[
 \Q_{\ell,X}\longrightarrow R\pi_*\Q_{\ell,X'}
 \xrightarrow{\operatorname{Tr}_\pi}\Q_{\ell,X}
\]
whose composite is multiplication by the generic degree on each component. Here the trace is not an unjustified trace for a nonfinite map: a projective generically finite map between regular schemes is a smoothable local complete-intersection morphism of virtual dimension zero. Absolute purity identifies its Gysin map with $\Q_\ell\to\pi^!\Q_\ell$, and proper adjunction gives the displayed trace. The degree computation is checked on the generic finite-flat locus, since an endomorphism of the constant sheaf on a connected scheme is a scalar. The precise construction and calculation are \citet[Corollary 3.1.2, Lemmas 3.5.8--3.5.10, and Proposition 2.5.13]{c_riou2014}; the equality also covers inseparable generic degree by the finite-flat degree formula. Dividing by the degree proves that $Rf_*\Q_\ell$ is a direct summand of $Rh_*Rv_*\Q_\ell$.

Let $D^{(q)}$ denote the disjoint union of the intersections of $q$ distinct boundary components, and put $D^{(0)}=\widetilde X$. Absolute purity gives
\[
 R^qv_*\Q_\ell=\Q_{\ell,D^{(q)}}(-q).
\]
This is \citet[Corollary 3.1.4]{c_riou2014}, equivalently \citet[Theorem 2.1.2]{c_pillonistroh2014}; it remains true for regular arithmetic schemes. The stalkwise Leray sequence and proper base change now give
\[
 E_2^{a,q}=H^a(F_q,\Q_\ell)(-q)
 \Longrightarrow (R^{a+q}(f\pi)_*\Q_\ell)_{\bar y},
 \qquad F_q=D^{(q)}\times_Y\bar y.
\]
Every $F_q$ is proper. By the definition of $n$ and regular codimension, $\dim F_q\leq n-q$ for $q\geq1$, while $\dim F_0\leq n-\epsilon$. This step does not assume that intersecting a special fibre with a vertical divisor lowers its dimension: the inequality already holds for the entire local scheme $D^{(q)}$, of codimension $q$ in a regular local source of dimension at most $n$. The extra drop for $F_0$ follows because a proper closed fibre containing a component of maximal local dimension would make that component map to $y$. Alteration does not change the image of its generic point.

Write $a=i-q$. Lemma~\ref{c:proper-level} says that each constituent of $E_2^{a,q}$ is $V_r(-(t+q))$, with
\[
 r+2t\leq i-q,
 \qquad t\geq\max(0,i-q-\dim F_q).
\]
For $q=0$ this gives $t\geq L$. For $q\geq1$ it gives $t+q\geq q+\max(0,i-n)\geq L$. Its total weight $w=r+2(t+q)$ satisfies $w\leq i+q$, and nonvanishing requires $q\leq\min(i,2n-i)$ (as well as $q\leq n$). It follows that $w\leq2\min(i,n)$ and $r\leq w-2L$. The other inequality $r\leq2\min(i,n)-w$ follows as follows. If $i\leq n$, use $w\leq i+q$ to obtain $r=w-2(t+q)\leq2i-w$. If $i>n$, use $t\geq i-n$ for $q\geq1$, and combine $q\geq w-i$ with the same formula, obtaining $r\leq2n-w$; the $q=0$ term obeys it because $w\leq i$ and $t\geq i-n+\epsilon$. Thus every term lies in \eqref{c:weight-level-polygon}. A spectral sequence gives a filtration by subquotients of its terms; the inequalities survive extensions and direct summands. No degeneration of this Leray sequence is needed.

For the coefficient sheaf, relative hard Lefschetz and \citet[Theorem 1.5 and Proposition 2.4]{c_deligne1968} give a derived splitting
\[
 Rg_*\Q_\ell\simeq\bigoplus_bR^bg_*\Q_\ell[-b].
\]
Accordingly $R^if_*F$ is a direct summand of $R^{i+j}(fg)_*\Q_\ell$. The source of $g$ is regular, its local dimension is $n+j$, and the value of $\epsilon$ is unchanged. Applying the constant-coefficient result with $(i,n)$ replaced by $(i+j,n+j)$ gives exactly \eqref{c:weight-level-polygon} and the displayed bounds. This proves the corrected statement, including the margin condition.
\end{proof}

\begin{proposition}[Optimality and the limit of the word sharp]
\label{c:level-optimality}
\status{Proved optimality in the stated senses}
For constant coefficients the preceding upper weight and level bounds are attained for every nonnegative degree in the indicated nonzero ranges. The general and margin level bounds remain optimal after adding $j$. The coefficient $2j$ in the universal weight envelope cannot be replaced by $j$: for every $j\geq0$ there is a strictly special weight-$j$ example with $n=i=1$ whose $R^1$ boundary stalk is pure of weight $2j+2$.
\end{proposition}

\begin{proof}
For the general level bound, take the structural map of $E^n$ over $\F_5$, where $E$ is the ordinary elliptic curve of Example~\ref{c:finite-field-lie}. Its Frobenius eigenvalues have $5$-adic slopes $0,1$. In $H^i(E^n)$ the smallest slope is $\max(0,i-n)$ and the weight is $i$. Any special constituent $V_r(-t)$ has nonnegative slopes before twisting, since eigenvalues on smooth projective cohomology are algebraic integers. Hence such a lowest-slope eigenvalue requires $t\leq\max(0,i-n)$, and therefore $r=i-2t\geq\min(i,2n-i)$. The upper bound is attained. Tensoring with the constant strictly special object $H^j(E^j)$, which has a slope-zero eigenvalue, raises the required level by exactly $j$.

For the margin case let $B=E^{n-1}$ and choose a sufficiently ample projectively normal embedding. Let $Y$ be its affine cone and $X=Y\setminus\{y\}$, where $y$ is the vertex. Then $X$ is the complement of the zero section of a line bundle on the smooth $B$, so it is regular of dimension $n$ and $X\hookrightarrow Y$ is an open immersion. The blowup of the vertex gives the line-bundle resolution. Purity and proper base change identify its local cohomology with the Gysin sequence for that line bundle:
\[
 \operatorname{coker}\bigl(L:H^{i-2}(B)(-1)\to H^i(B)\bigr)
 \longrightarrow (R^if_*\Q_\ell)_{\bar y}
 \longrightarrow\ker\bigl(L:H^{i-1}(B)(-1)\to H^{i+1}(B)\bigr).
\]
The two outer terms give a short exact sequence. For $i\leq n-1$, a slope-zero class in $H^i(B)$ cannot be in the image of $L$, whose eigenvalues have slopes at least one. It gives level $i$. For $n\leq i\leq2n-1$, put $r=2n-i-1$. The primitive degree-$r$ cohomology contains a slope-zero class for the same reason. Hard Lefschetz carries it to the right-hand kernel after $i-n$ powers of $L$, and the boundary twist adds one further Tate twist. This has weight $i+1$ and slope $i-n+1$, forcing level at least $2n-i-1$. Tensoring again with $H^j(E^j)$ proves sharpness with $+j$. For $n=1$ the same construction is simply $\Gm\hookrightarrow\mathbb A^1$.

For weights with $j=0$, if $i\leq n$ use the origin-boundary stalk of
$\Gm^i\times\mathbb A^{n-i}\hookrightarrow\mathbb A^n$, whose degree-$i$ Kummer class has weight $2i$. If $n<i\leq2n$, put $q=2n-i$ and $r=i-n$. The product of the degree-$q$ class of $\Gm^q$ with the top class of $\mathbb P^r$ has degree $i$ and weight $2(q+r)=2n$. Take the product open immersion followed by projection to $\mathbb A^q$; if $q=0$ this is the projective structural map. For the margin range $i\leq2n-1$ we have $q\geq1$, so the same example satisfies the extra condition.

Finally let $U=\mathbb A^1_{\F_5}\setminus\{0,1\}$ and let $\pi:\mathcal E\to U$ be the Legendre family
\[
 y^2=x(x-1)(x-t).
\]
At $t=0$ its discriminant $16t^2(1-t)^2$ has order two and $c_4=16(1-t+t^2)$ is a unit. The node has tangent slopes whose squares are $-1$, which is a square in $\F_5$, so the reduction is split multiplicative. Thus on $H^1$ inertia acts by a nonzero rank-one unipotent operator, with invariant line $\Q_\ell$ and coinvariant line $\Q_\ell(-1)$. This is the Tate-curve calculation underlying SGA~7~I, IX, Proposition 2.2.5. Take $F_j=\operatorname{Sym}^j R^1\pi_*\Q_\ell$. The K\"unneth and Schur idempotents exhibit it as a direct summand of $R^j(\mathcal E^j/U)_*\Q_\ell$, so it has the required dimension-$j$ presentation. Its inertia coinvariants are $\Q_\ell(-j)$, and tame cohomology supplies one further twist:
\[
 (R^1u_*F_j)_{\overline0}=\Q_\ell(-j-1),
 \qquad u:U\hookrightarrow\mathbb A^1.
\]
This proves the claimed weight $2j+2$. Products with the Kummer and projective examples above attain $2\min(i,n)+2j$ for every $1\leq i\leq2n-1$. The finite matrix calculation of the symmetric-power monodromy and the displayed optimization inequalities are recorded in the accompanying program.
\end{proof}

The theorem proves the requested upper-weight inequality, the strongest displayed uniform geometric-level bounds, and the manuscript's margin question in the explicit finite-type setting. We do not silently call the weight envelope pointwise optimal at the remaining coefficient endpoints. Already for an open arithmetic trait, $i=0$ and $F=H^j$ of a smooth projective generic fibre, improving the bound $2j$ to $j$ says that inertia invariants have weights at most $j$; this is the local-invariant-cycle problem. We retain that sharper endpoint as a separate gate rather than derive it from a theorem about ordinary direct-image weights. The same caution applies if ``strictly special'' is defined without the dimension-$j$ clause of 12:111: the theorem above proves exactly the dimension-normalized convention, and a relative dimension-reduction argument must be supplied for the broader convention.
